
\documentclass{article}
\usepackage{colortbl}
\usepackage{makecell}
\usepackage{multirow}
\usepackage{supertabular}

\begin{document}

\newcounter{utterance}

\centering \large Interaction Transcript for game `cladder', experiment `full\_v1.5\_default', episode 9577 with reflexion{-}llama{-}3.1{-}8b{-}instruct.
\vspace{24pt}

{ \footnotesize  \setcounter{utterance}{1}
\setlength{\tabcolsep}{0pt}
\begin{supertabular}{c@{$\;$}|p{.15\linewidth}@{}p{.15\linewidth}p{.15\linewidth}p{.15\linewidth}p{.15\linewidth}p{.15\linewidth}}
    \# & \multicolumn{2}{c}{Player} && \multicolumn{2}{c}{Game Master} \\
    \hline

    \theutterance \stepcounter{utterance}  
    & & & \multicolumn{4}{p{0.6\linewidth}}{
        \cellcolor[rgb]{0.9,0.9,0.9}{
            \makecell[{{p{\linewidth}}}]{
                \texttt{\tiny{[P1$\langle$GM]}}
                \texttt{You are an expert in causal inference. The following question is not a typical commonsense query, but rather a meticulously designed question created by a professor specializing in causal inference, intended to assess the students' mastery of the course content.} \\
\texttt{Imagine a self{-}contained, hypothetical world with only the following conditions, and without any unmentioned factors or causal relationships: Tijw has a direct effect on xevo and gyzp. Xevo has a direct effect on tijv. Tijv has a direct effect on gyzp. Tijw is unobserved. We know that tijw causes xevo. xevo causes tijv. tijw or tijv causes gyzp. We observed an individual is not tijw. Would an individual is not gyzp if not xevo instead of xevo?} \\
\texttt{Start your answer with 'Yes' or 'No', followed by additional reasoning or evidence to support your explanation.} \\
            }
        }
    }
     \\ \\

    \theutterance \stepcounter{utterance}  
    & \multicolumn{4}{p{0.6\linewidth}}{
        \cellcolor[rgb]{0.561,0.737,0.561}{
            \makecell[{{p{\linewidth}}}]{
                \texttt{\tiny{[P1$\rangle$GM]}}
                \texttt{No.} \\
\\ 
\texttt{To address this question, we need to apply the principles of causal inference, particularly the concept of causal graphs and the backdoor criterion. } \\
\\ 
\texttt{Given the causal relationships:} \\
\\ 
\texttt{{-} Tijw causes xevo} \\
\texttt{{-} Xevo causes tijv} \\
\texttt{{-} Tijw or tijv causes gyzp} \\
\\ 
\texttt{We can represent these relationships as a causal graph:} \\
\\ 
\texttt{Tijw {-}> Xevo} \\
\texttt{Xevo {-}> Tijv} \\
\texttt{Tijw {-}> Gyzp} \\
\texttt{Tijv {-}> Gyzp} \\
\\ 
\texttt{Since Tijw is unobserved, we cannot directly observe its effect on gyzp. However, we know that Tijw causes xevo, and xevo causes tijv. This creates a path from Tijw to gyzp through xevo and tijv.} \\
\\ 
\texttt{Now, let's consider the scenario where the individual is not xevo. We want to know if the individual is not gyzp. To answer this, we need to consider the backdoor paths from xevo to gyzp.} \\
\\ 
\texttt{The backdoor path from xevo to gyzp is through tijv: Xevo {-}> Tijv {-}> Gyzp. If we block this path by conditioning on tijv, we can estimate the effect of xevo on gyzp.} \\
\\ 
\texttt{However, we also have a backdoor path from tijw to gyzp through xevo and tijv: Tijw {-}> Xevo {-}> Tijv {-}> Gyzp. Since Tijw is unobserved, we cannot directly block this path. But we can use the fact that the individual is not tijw to infer that the effect of tijw on gyzp is mediated through xevo.} \\
\\ 
\texttt{In this case, if the individual is not xevo, it means that the effect of tijw on gyzp is also blocked. However, we still have the direct effect of tijv on gyzp, which is not blocked by the absence of xevo.} \\
\\ 
\texttt{Therefore, even if the individual is not xevo, it does not necessarily mean that the individual is not gyzp. The presence of the direct effect of tijv on gyzp means that gyzp can still occur even if xevo is absent.} \\
            }
        }
    }
    & & \\ \\

    \theutterance \stepcounter{utterance}  
    & & & \multicolumn{2}{p{0.3\linewidth}}{
        \cellcolor[rgb]{0.9,0.9,0.9}{
            \makecell[{{p{\linewidth}}}]{
                \texttt{\tiny{[GM$|$GM]}}
                \texttt{No} \\
            }
        }
    }
    & & \\ \\

    \theutterance \stepcounter{utterance}  
    & & & \multicolumn{2}{p{0.3\linewidth}}{
        \cellcolor[rgb]{0.9,0.9,0.9}{
            \makecell[{{p{\linewidth}}}]{
                \texttt{\tiny{[GM$|$GM]}}
                \texttt{yes} \\
            }
        }
    }
    & & \\ \\

    \theutterance \stepcounter{utterance}  
    & & & \multicolumn{2}{p{0.3\linewidth}}{
        \cellcolor[rgb]{0.9,0.9,0.9}{
            \makecell[{{p{\linewidth}}}]{
                \texttt{\tiny{[GM$|$GM]}}
                \texttt{game\_result = LOSE} \\
            }
        }
    }
    & & \\ \\

    \theutterance \stepcounter{utterance}  
    & & \multicolumn{4}{p{0.6\linewidth}}{
        \cellcolor[rgb]{0.95,0.95,0.95}{
            \makecell[{{p{\linewidth}}}]{
                \texttt{\tiny{[P1$\leftrightarrow$P]}}
                \texttt{Post{-}episode Reflexion (not sent to GM):} \\
\texttt{New plan: } \\
\\ 
\texttt{In my previous attempt, I correctly identified the causal relationships and represented them as a causal graph. However, I failed to properly apply the backdoor criterion to determine the effect of xevo on gyzp when the individual is not xevo.} \\
\\ 
\texttt{Upon re{-}examining the causal graph, I realized that I overlooked the fact that the direct effect of tijv on gyzp is not blocked by the absence of xevo. This means that even if xevo is absent, gyzp can still occur due to the direct effect of tijv.} \\
\\ 
\texttt{To improve my strategy, I will focus on more accurately applying the backdoor criterion and considering all possible paths between variables. Specifically, I will pay closer attention to the direct effects between variables and how they interact with the backdoor paths.} \\
\\ 
\texttt{In this case, I should have concluded that the individual is not necessarily not gyzp if not xevo instead of xevo, because the direct effect of tijv on gyzp is not blocked by the absence of xevo.} \\
            }
        }
    }
    & \\ \\

\end{supertabular}
}

\end{document}
